\section{The construction problem and the order of its limits}
\label{ex-statement}

The aim is to construct a relativistic field theory from finite-dimensional
Wilson integrals. There are two limiting operations. First the lattice spacing
tends to zero while the physical period stays fixed. Then the period tends to
infinity. The observables and bare action coefficients must be chosen coherently
before either operation: changing the family in response to a requested sequence
of periods would give a weaker statement.

\paragraph{A caution about the target of this section.}
The target formulated below is the statement that the source verifies by
machine. Read on its own, it is weaker than a reader might expect from the
words ``existence of four-dimensional Yang--Mills theory''. It can be
satisfied by a generalized free scalar field, that is, a quantum field whose
vacuum correlations are determined by its two-point function through sums
over pairings, here with a continuous distribution of masses; and the lattice
observables that the target asks for can then be chosen independently of the
gauge configuration. Section~\ref{ex-target-free-witness} states this
precisely and explains which features of the target allow it. The target
alone therefore does not express the specifically Yang--Mills content of the
construction. That content lies in how the later sections reach the target:
they work with specific local gauge-invariant functions of the Wilson links,
and every limiting object is obtained from the Wilson measure. The
observation does not affect that construction. It does mean that a main
theorem intended to express this content should also record the chosen
lattice observables and their connections to the limiting theory. The source
has such a detailed statement, and the group reviewing the source bases its
assessment on a stronger target of this kind
(Section~\ref{ex-target-free-witness}); formulating it as the main theorem of
this article remains to be done.

We assume familiarity with finite-dimensional Lie groups, probability Haar
measure on compact groups, Lebesgue integration and dominated convergence,
and elementary Hilbert-space theory. The specific compactness, quotient and
operator arguments used below are supplied where they enter the construction.

Fix a closed subgroup \(G\subset U(r_G)\), with its specified faithful matrix
representation. Its real Lie algebra is
\[
 \mathfrak g=\{X\in M_{r_G}(\mathbb C):\exp(tX)\in G
                       \text{ for every }t\in\mathbb R\}.
\]
We assume that \(\mathfrak g\) is nonabelian and simple: its only real
subspaces invariant under every \(\operatorname{ad}X\), \(X\in\mathfrak g\),
are \(0\) and \(\mathfrak g\). The group \(G\) is compact, being closed in
\(U(r_G)\). It is not assumed to be connected, and the given representation
is not assumed to be irreducible. The Wilson action of
Section~\ref{ex-lattice} is normalized by the dimension \(r_G\) of this
representation, so the same abstract group in two representations gives two
instances of the problem. Coordinates on \(\mathbb R^4\) are written
\(x=(\boldsymbol x,t)\), with time last. The Minkowski form is
\(t^2-|\boldsymbol x|^2\), and the closed forward cone is
\(\overline V_+=\{(\boldsymbol p,E):|\boldsymbol p|\le E\}\).

The lattice family consists of strictly increasing positive integer sequences
\(P_p\) and \(D_N\), a spacing \(a_N=D_N^{-1}\), real coefficients
\(\beta_{N,p}\), and local gauge-invariant insertion functions
\(O_{N,p,v}(x;U)\) indexed by tokens \(v\). The periodic lattice has
\(P_pD_N\) sites along each axis and physical period \(P_p\).
Section~\ref{ex-lattice} describes its Wilson probability measure and moments.
We call the family \emph{regular} when its insertions are continuous in
\(U\), invariant under every lattice gauge transformation, and uniformly
local in the following precise sense. For each token \(v\) there is a finite
stencil \(\Sigma_v\subset\mathbb Z^4\), independent of \(N,p,x\), such
that two configurations agreeing on all links with both endpoints in
\(x+\Sigma_v\), reduced modulo \(P_pD_N\), give the same insertion value.
The fourth regularity condition is
\begin{equation}\label{ex-running-input}
 \exists\sigma>0\quad\forall\varepsilon>0\quad\exists N_0\quad
 \forall N\ge N_0\quad\forall p:\qquad
 \left|\frac{\beta_{N,p}}{\log(a_N^{-1})}-\sigma\right|<\varepsilon.
\end{equation}
In particular, the choice of \(N_0\) is uniform in the period index.

Let \(\operatorname{Conf}_n(\mathbb R^4)\) be the set of ordered
\(n\)-tuples of pairwise distinct points. A separated test is a smooth
complex function with compact support in this set, extended by zero outside
it. For insertion tokens \(\boldsymbol v=(v_1,\ldots,v_n)\), write
\(M_{N,p}(\boldsymbol v;\boldsymbol x)\) for their Wilson moment and put
\begin{equation}\label{ex-smearing}
 W_{N,p}(\boldsymbol v;f)=a_N^{4n}
 \sum_{\boldsymbol x\in\Lambda_{N,p}^n}
 f(a_N\boldsymbol x)M_{N,p}(\boldsymbol v;\boldsymbol x).
\end{equation}
Here the coordinates of \(\boldsymbol x\) are chosen in the centered
fundamental cell. A period is suitable for \(f\) if its support lies strictly
inside that physical cell in every coordinate.

The limiting data have several related components. For every integer grade
\(d\ge0\) there is a finite-dimensional complex space \(\mathcal F_d\).
It has a complex-linear injection into a common field-label space; every label is
represented at some grade. Tokens map onto the disjoint union
\(\coprod_d\mathcal F_d\). A compact metric space \(X\), a probability
measure \(\nu\) of full support, and a selected character \(\chi_*\in X\)
describe the phases. For graded fields \(\boldsymbol A\), let
\(S^{\chi}_{\boldsymbol A}(\boldsymbol x)\) and
\(S^{\mathrm{mix}}_{\boldsymbol A}(\boldsymbol x)\) denote the phase and
mixed separated correlations. They are smooth on configuration space,
permutation symmetric, and continuous in \(\chi\), with
\begin{equation}\label{ex-mixture}
 S^{\mathrm{mix}}_{\boldsymbol A}(\boldsymbol x)
   =\int_X S^{\chi}_{\boldsymbol A}(\boldsymbol x)\,d\nu(\chi).
\end{equation}
The mixed correlations are the Wilson limits. The selected correlations
are the Euclidean correlations to be continued into the reconstructed theory.
These two roles should be kept distinct even when the phase space has only
one point. At the level of this target, only the selected point is required
to carry a reconstructed relativistic theory. The other points are required
to satisfy the displayed correlation properties and the common
operator-product estimates. Multilinearity of all these correlations in the
graded fields is not an additional condition of the target.

\begin{definition}[Reconstructed selected theory]\label{ex-reconstructed}
The selected theory consists of a complex Hilbert space \(\mathcal H\),
a normalized vacuum \(\Omega\), a common dense invariant domain
\(\mathcal D\) containing \(\Omega\), and complex-linear fields
\(\Phi(A,f):\mathcal D\to\mathcal D\), linear separately in the label
and test, where
\(f\in\mathcal S(\mathbb R^4)\). For fixed \(A,\psi\), the map
\(f\mapsto\Phi(A,f)\psi\) is continuous in the Schwartz topology.
The domain is the span of field words applied to the vacuum. There is a
conjugate-linear label adjoint \(A\mapsto A^*\) with
\[
 \langle\Phi(A,f)\psi,\eta\rangle
   =\langle\psi,\Phi(A^*,\overline f)\eta\rangle.
\]
A strongly continuous unitary representation of the proper orthochronous
Poincar\'e group preserves the domain and the vacuum, and transforms fields
with the specified action on labels and the ordinary affine pullback of tests.
Its translation spectral measure is supported in \(\overline V_+\), and
every translation-invariant vector is a scalar multiple of \(\Omega\).
Fields commute on \(\mathcal D\) when their test supports are spacelike
separated. Every nonzero label acts nontrivially for some test, and at least
one smeared field is nonzero.

For fixed spatial positions the selected correlations have a holomorphic
continuation \(F_{\boldsymbol A}(\boldsymbol z;\boldsymbol x)\) on
\(\operatorname{Im}z_1<\cdots<\operatorname{Im}z_n\). At strictly ordered
Euclidean times its value at \(z_i=it_i\) is
\(S^{\chi_*}_{\boldsymbol A}((\boldsymbol x_i,t_i)_i)\).
For every strictly increasing real height vector \(\boldsymbol h\) and
all \(f_1,\ldots,f_n\in\mathcal S(\mathbb R^4)\),
\begin{equation}\label{eq:public-selected-boundary}
\begin{split}
 \lim_{\epsilon\downarrow0}\int_{(\mathbb R^4)^n}
 &F_{\boldsymbol A}\bigl((t_j+i\epsilon h_j)_j;
                  (\boldsymbol x_j)_j\bigr)
       \prod_{j=1}^n f_j(\boldsymbol x_j,t_j)\,
                         d\boldsymbol x\,d\boldsymbol t\\
 &=\langle\Omega,\Phi(A_1,f_1)\cdots\Phi(A_n,f_n)\Omega\rangle.
\end{split}
\end{equation}
In this target clause the integral has value zero when its integrand is not
Lebesgue integrable. Thus the clause prescribes a boundary limit without
separately requiring absolute integrability at positive height. When ordinary
integrals are used in the construction, their integrability must be supplied
by the analytic estimates, as in Lemma~\ref{lem:tested-cr} and
Theorem~\ref{thm:joint-boundary}.
\end{definition}

The existence target is one regular Wilson family such that, for every
strictly increasing period-index sequence \(p(k)\), a strictly increasing
\(q(k)\) and one hierarchy with the preceding properties can be found.
For each separated test there are row limits \(w_p(\boldsymbol v;f)\)
satisfying
\begin{align}
 \lim_{N\to\infty}W_{N,p}(\boldsymbol v;f)
   &=w_p(\boldsymbol v;f)
       &&\text{at each suitable fixed }p,\label{ex-row-target}\\
 \lim_{k\to\infty}w_{p(q(k))}(\boldsymbol v;f)
   &=\int f(\boldsymbol x)
           S^{\mathrm{mix}}_{\boldsymbol A}(\boldsymbol x)\,d\boldsymbol x,
       \label{ex-volume-target}
\end{align}
where \(A_i\) is the class of \(v_i\). Suitability holds eventually
along the selected period sequence. The same \(q\) and hierarchy work for
all arities, tokens and tests. The hierarchy may depend on the requested
sequence; the Wilson family may not.

The full endpoint also specifies a common operator-product expansion and a
short-distance test on one centered field of grade four, described in
Section~\ref{ex-assembly}. These are additional obligations on the same
hierarchy. The scope of the target should be read literally. The spectral
conclusions are that the energy-momentum spectrum lies in the closed forward
cone and that the translation-invariant vectors form the line
\(\mathbb C\Omega\); no gap above the vacuum in the spectrum is asserted.
The short-distance test prescribes the leading behaviour
\(r^{-8}\log^{-2}(r^{-1})\) of one two-point function on the time axis and,
as a consequence, excludes its coincidence near \(r=0\) with any finite
combination of the mass-shell profiles of Section~\ref{sec:trace-profiles};
it does not exclude infinite or continuous superpositions of masses and makes
no assertion about truncated correlations of higher order. The target
requires a positive limit of \(\beta_{N,p}/\log(a_N^{-1})\) but does not
prescribe its value. The grade-four field is existentially chosen, subject
to the stated test (vanishing vacuum one-point functional, the exact
asymptotic coefficient, and the profile exclusion); the target contains no
stress-tensor or Ward-identity normalization identifying it physically, and
that identification belongs to the detailed construction. More generally the
target concerns a theory of the constructed gauge-invariant observables: the
insertions \(O_{N,p,v}\) are part of what is constructed, no lattice formula
is prescribed for them beyond regularity and the joint convergence
requirements, and the lattice measure enters the target only through the
mixed correlations. Identifying a particular label with a separately
prescribed physical observable requires its own construction argument.
The detailed construction is parameterized by a positive blocking penalty;
the target is obtained from the value \(1\) in the conditional assembly of
Proposition~\ref{ex-conditional-existence}. No uniqueness assertion is made.

\begin{remark}[What the target alone ties to the Wilson measure]
\label{rem:ex-target-adequacy}
The Wilson measure enters the target only through
\eqref{ex-row-target}--\eqref{ex-volume-target}, that is, through the mixed
correlations. These are tied to the selected correlations only by the
mixture identity \eqref{ex-mixture} and by the common operator-product
estimates. The target does not require the correlations to be multilinear in
the graded fields, nor the correlations at a non-selected phase to be
reflection positive. Consequently the public clauses alone do not force the
selected theory to arise from the Wilson measure: they admit Wilson families
whose insertions do not depend on the gauge configuration.

To see this, suppose hierarchy data are given that satisfy every requirement
of the target except the two convergence statements, with selected
correlations \(S=S^{\chi_*}\), and suppose that some \(u\in\mathcal F_0\),
\(u\ne0\), enters \(S\) homogeneously: replacing the tag \((0,u)\) by
\((0,2u)\) in any one slot doubles \(S\). Put \(m(A)=1\) if \(A=(0,u)\) and
\(m(A)=0\) for every other tag, and \(M_{\boldsymbol A}=\prod_im(A_i)\), with
\(M_\emptyset=1\). Let the tokens be the tags themselves, with the identity
token map; take \(P_p=p+1\), \(D_N=N+2\), \(\beta_{N,p}=\log(N+2)\) and the
insertions \(O_{N,p,A}(x;U)=m(A)\), which are constant, hence regular with
empty stencils, and satisfy \eqref{ex-running-input} with \(\sigma=1\).
Since the Wilson measure is a probability measure, the moments are
\(M_{\boldsymbol A}\), and Riemann sums give
\(W_{N,p}(\boldsymbol A;f)\to M_{\boldsymbol A}\int f\) at every suitable
\(p\). Replace the phase data by the two-point space \(\{+,-\}\) with the
uniform measure and selected point \(+\), and set
\[
 S^+=S,\qquad S^-=2M-S,\qquad S^{\mathrm{mix}}=M .
\]
The mixture identity holds, the selected theory is unchanged, and the
volume limit is \(\int fM\) along every period sequence. In the
operator-product data add two grade-zero channels with actors \(u\) and
\(2u\) and coefficients \(h_a\) and \(-\tfrac12h_a\), where
\(h_a=\prod_km(a_k)-\sum_ic^{0,i}_a\,m((0,E_{0,i}))\) is the defect of the
old grade-zero channels \((E_{0,i},c^{0,i}_a)\) on \(M\). The coefficient
bounds for the new channels follow from those of the old ones, because
\(|m|\le1\). At the selected point the two new terms cancel by homogeneity;
in the mixed state the expansion reproduces \(M\) exactly, because \(M\)
vanishes on every actor of positive grade; and the remainder at \(-\) is
minus that at \(+\). Every requirement other than the convergence statements
mentions only the selected point, the common operator-product data or the
listed properties of the phase correlations, and these have just been
checked. The lattice moment of every field of positive grade, including the
grade-four field of the trace test, is zero in this example.

The detailed construction does not use any of these freedoms: its
correlations are multilinear in the fields, each of its phases carries its
own positive reconstruction, and its insertions are fixed local functions of
the links tied to the Wilson measure by the construction. The remark
therefore concerns only what the displayed target certifies; it says nothing
against the construction of the later sections.
\end{remark}

\subsection{The target is satisfied by a generalized free field}
\label{ex-target-free-witness}

Remark~\ref{rem:ex-target-adequacy} supplies the lattice clauses of the
target, with insertions that do not depend on the gauge configuration, once
hierarchy data are given that satisfy the remaining requirements and the
grade-zero homogeneity hypothesis of the remark. Such data need not come from
a gauge theory.

\begin{proposition}[A generalized free-field solution of the target]
\label{prop:ex-target-free-witness}
For every group \(G\) as in this section, the existence target, including the
operator-product requirements
\eqref{ex-ope-coefficient}--\eqref{ex-ope-remainder} and the trace
requirement \eqref{ex-trace-limit} with its profile exclusion, is satisfied by
a regular lattice family whose insertions do not depend on \(U\) and by a
hierarchy whose selected theory is the field system of Wick polynomials and
derivatives of one generalized free scalar field. The continuum data depend
on \(G\) only through \(\dim_{\mathbb R}\mathfrak g\) in the normalization of
the grade-four field of the trace requirement.
\end{proposition}

The proposition has been verified by machine against the formal statement of
the source, using that statement and none of the source's construction. The
record, with the formal development and its reviews, is in the blueprint
(unit YM-000026, result G01). We describe the construction and do not repeat
the proofs; the formal proofs have been checked by machine but not yet read
by a mathematician.

\paragraph{The field.}
Let \(\mu\) be the measure on \(\mathbb R^4\) with density
\[
 w(p)=\frac{(p^2)^2}{\log^2(e+p^2)},\qquad p^2=E^2-|\boldsymbol p|^2,
\]
on the closed forward cone \(\overline V_+\) and zero outside it. It is
invariant under the proper orthochronous Lorentz group and polynomially
bounded. Let \(\phi\) be the generalized free scalar field with two-point
distribution \(\int e^{i\langle y-x,p\rangle}\,d\mu(p)\), where
\(\langle\cdot,\cdot\rangle\) is the Euclidean pairing of \(\mathbb R^4\) and
no normalizing factor is inserted. Relative to the invariant measure
\(d^3\boldsymbol p/(2\sqrt{s+|\boldsymbol p|^2})\) on the mass shell of mass
squared \(s\), its K\"all\'en--Lehmann measure is
\(s^2\log^{-2}(e+s)\,ds\) on \(s\ge0\): the mass spectrum is continuous. Write
\(\xi_\alpha=i^{\alpha_t}\partial^\alpha\phi\) for a multi-index
\(\alpha\in\mathbb N_0^4\) with time component \(\alpha_t\); the factor \(i\)
for each time derivative matches the Euclidean and relativistic conventions.
The labels are the finite complex linear combinations of Wick monomials
\(:\!\xi_{\alpha_1}\cdots\xi_{\alpha_k}\!:\). A monomial has grade
\(\sum_j(4+|\alpha_j|)\), the empty monomial is the unit, and
\(\mathcal F_d\) is the span of the monomials of grade at most \(d\), which
is finite-dimensional. The density \(w\) is strictly positive in the interior
of the forward cone. Together with the linear independence of the
polynomials in the momenta that distinct Wick monomials produce, this shows
that every nonzero label acts nontrivially.

\paragraph{The relativistic theory.}
The Hilbert space is the closed span of the vectors obtained by applying field
words to the vacuum; these vectors are written down explicitly as families of
square-integrable functions of \(k\) momenta, \(k\ge0\). The properties of
Definition~\ref{ex-reconstructed} are verified directly: linearity, the
adjoint, continuity in the test, Poincar\'e covariance, the support of the
spectral measure in \(\overline V_+\), uniqueness of the vacuum (in every
sector with \(k\ge1\) the total energy is positive almost everywhere; no mass
gap is needed), and locality. Let
\[
 S(\xi)=\int e^{-|\xi|E}\,d\mu(p),\qquad \xi\in\mathbb R^4\setminus\{0\},
\]
which for \(\xi=(\boldsymbol u,\tau)\) with \(\tau>0\) equals
\(\int e^{-\tau E+i\boldsymbol u\cdot\boldsymbol p}\,d\mu(p)\). At pairwise
distinct insertion points, the selected correlations are sums over complete
pairings of the individual factors \(\xi_\alpha\), with no pair joining two
factors of the same Wick monomial; each pair contributes the corresponding
derivative of \(S\), and the correlation vanishes if there is no such
pairing. Their continuation to complex times and the boundary values
\eqref{eq:public-selected-boundary} are obtained from the same momentum
integrals.

\paragraph{Operator products.}
The expansion of a cluster of Wick monomials is given by Wick's theorem
followed by Taylor expansion, about the base point, of the factors that are
not contracted within the cluster. Each coefficient is a finite sum, over
partial contractions and Taylor multi-indices, of products of derivatives of
\(S\) at differences of cluster points with monomials in the displacements
and inverse factorials. One family of coefficients serves every truncation
degree. The Wick decomposition and the Taylor remainder, together with
\(|\partial^\gamma S(\xi)|\le C_\gamma|\xi|^{-8-|\gamma|}\), give the powers of
\(r\) required in \eqref{ex-ope-coefficient} and \eqref{ex-ope-remainder},
at first with constants depending on the fields, the truncation degree and
the geometry parameters. The separation of the spectators makes these
estimates uniform in the allowed positions, and for fixed parameters the
factor \(r^{-\epsilon}\) absorbs the constants.

\paragraph{The trace requirement.}
On the time axis, with \(e_t=(\boldsymbol0,1)\),
\[
 \lim_{r\downarrow0}r^8\log^2(r^{-1})\,S(re_t)=384\pi .
\]
Take \(\Theta=\lambda\phi\) with \(\lambda>0\) and
\(\lambda^2=\dim_{\mathbb R}\mathfrak g/(128\pi^5)\). Then \(\Theta\) has grade
four, its vacuum one-point functional vanishes, and \eqref{ex-trace-limit}
holds with the required constant. The profile exclusion follows from this
asymptotic alone, by Theorem~\ref{thm:profile-exclusion}.

\paragraph{The lattice clauses.}
The correlations just described are multilinear in the fields, so the unit
\(u\in\mathcal F_0\) enters them homogeneously, and
Remark~\ref{rem:ex-target-adequacy} applies: the tokens are the tags, the
insertions are the constants \(m(A)\), the phase space has two points, and
two grade-zero channels are added to the operator-product data. This
completes the data. The formal proof uses a variant of the remark in which
\(m\) is the indicator of the zero vector of \(\mathcal F_0\) and one
channel with zero field is added in each grade; it relies on the vanishing
of the correlations when a field is zero.

\paragraph{Why the target admits this.}
Three features of the target, each already noted above, combine.
\begin{enumerate}
\item The Wilson measure enters only through the mixed correlations, and
these are tied to the selected theory only by the mixture identity and the
common operator-product estimates. The target imposes no multilinearity on
the phase and mixed correlation functions and no positivity or
reconstruction at a non-selected phase.
\item No particular nonconstant gauge-invariant observable is prescribed for
a field of the hierarchy. Beyond the lattice conditions, the group and its
representation enter the requirements only through
\(\dim_{\mathbb R}\mathfrak g\) in \eqref{ex-trace-limit}.
\item The profile exclusion constrains one two-point function. It excludes
the specified finite sums of products of mass-shell kernels, while the
continuous spectral weight chosen above gives exactly the required
short-distance behaviour. The operator-product requirements are upper
bounds, and they hold for this hierarchy of Wick polynomials. No clause
imposes a condition on truncated correlations of higher order that would
exclude a generalized free generating field.
\end{enumerate}

\paragraph{Relation to the review of the source.}
The weakness of the target is not a new observation in its general form. The
assessment of the source by its review group\footnote{Yang--Mills group,
\emph{Validation of the OpenAI Yang--Mills construction}, release of
10~November 2026, \url{https://github.com/yang-mills-group/YM}, Sections~7.3
and~9.} records that the authors of the source informed the group that the
original public target ``omitted some of the connections to the lattice
construction established by its proof'', and it bases its conclusions on a
stronger combined target that retains the lattice observables, limits, phases
and reconstructed fields of the construction itself. The same assessment
notes that the trace requirement excludes the specified finite free-field
descriptions ``but does not rule out every description in terms of composites
of Gaussian fields''. Proposition~\ref{prop:ex-target-free-witness} and
Remark~\ref{rem:ex-target-adequacy} make both remarks precise for the
original target: the omitted connections can be dispensed with altogether,
and the hierarchy of Wick polynomials and derivatives of one generalized free
scalar field supplies the remaining continuum requirements. They give no
information about the stronger combined target, which refers to the lattice
family of the construction.

\paragraph{What this does and does not say.}
The proposition concerns the target as a specification. It shows that the
target alone does not express the specifically Yang--Mills content of the
construction: a proof of it need not involve Yang--Mills theory. It does not
show any defect in the construction. The detailed theorem of the source
concerns mixed correlations obtained as limits of Wilson expectations of a
fixed family of local gauge-invariant insertions, a phase decomposition with
a positive reconstruction at every phase, and a grade-four field identified
with the centered trace of a stress tensor that satisfies translation and
rotation Ward identities; none of this is touched by the proposition. Nor
does the proposition bear on the question whether the theory so constructed
could itself be a generalized free field. That question is not settled by
the target, since the target is satisfied by one, and it is not addressed in
this article.
